\chapter{Linux简单查询手册与转行建议} \label{ch:linux}

嵌入式开发并不只有单片机。K230、RK3588、树莓派等平台都要运行 Linux，现代工具链（GCC、 CMake、Git、串口与网络调试）也大量依赖 Linux 环境。对只会 Keil 和 Windows 的开发者来说， Linux 第一次接触可能很陌生，但只要掌握文件、文本、进程、网络四类命令，再学会 GCC 与 CMake，就能立刻开始做嵌入式 Linux 开发。

\medskip 本章涉及到的实验内容：
\begin{enumerate}
\item Linux 虚拟机安装与基本命令；
\item GCC 单文件与多文件编译；
\item Makefile 自动构建；
\item CMake 工程构建；
\item 测试环节（由读者自行完成）。
\end{enumerate}

% ---------------------------------------------------------------------
% 以下为 Linux 命令全参数详解
本章下面将 Linux 常用命令参数进行全方位详解，涵盖文件操作、文本处理、进程管理、网络工具及 Shell 编程语法。所有示例均可在 Ubuntu 或 CentOS 下验证。

\section{文件与目录操作}

\subsection{\texttt{ls} —— 列出目录内容} 语法：\texttt{ls [选项] [路径]}

高频参数：
\begin{itemize}
\item \texttt{-l}：长格式（权限、链接数、所有者、组、大小、时间、文件名）。
\item \texttt{-a}：显示所有文件（含 \texttt{.} 开头的隐藏文件）。
\item \texttt{-A}：显示所有文件，但不显示 \texttt{.} 和 \texttt{..}。
\item \texttt{-h}：人类可读大小（1K、234M、2G）。
\item \texttt{-t}：按修改时间从新到旧排序。
\item \texttt{-r}：反转排序。
\item \texttt{-S}：按文件大小从大到小排序。
\item \texttt{-R}：递归子目录。
\item \texttt{-i}：显示 inode 号。
\item \texttt{-d}：仅列出目录本身（与 \texttt{-l} 联用）。
\item \texttt{-F}：在文件名后添加标识（\texttt{/} 目录，\texttt{*} 可执行，\texttt{@} 符号链接）。
\item \texttt{-1}：每行一个文件。
\item \texttt{--color=auto}：颜色区分文件类型。
\end{itemize}

技巧：
\begin{itemize}
\item \texttt{ls -lath}：按时间倒序、显示所有、人类可读、长格式。
\item \texttt{ls -ldh */}：查看当前目录下所有目录的大小。
\item \texttt{ls -lSrh}：按大小倒序，配合 \texttt{head} 找最大文件。
\end{itemize}

示例：
\begin{lstlisting}[language=bash]
	ls -lht /etc/*.conf
	ls -ld $(find . -maxdepth 1 -type d)
	ls -lt | head -11
\end{lstlisting}

注意：文件名含空格时建议使用 \texttt{--quoting-style=shell}。

\subsection{\texttt{cd} —— 切换目录} 语法：\texttt{cd [目标路径]}

特殊符号：
\begin{itemize}
\item \texttt{\string~}：家目录。
\item \texttt{-}：上一次所在目录（非常实用）。
\item \texttt{..}：上级目录。
\item \texttt{.}：当前目录。
\item \texttt{/}：根目录。
\end{itemize}

环境变量：\texttt{CDPATH} 可设置查找目录列表。

示例：
\begin{lstlisting}[language=bash]
	cd /var/log
	cd ../tmp
	cd -        # 回到 /var/log
	cd ~
\end{lstlisting}

注意：\texttt{cd} 是 shell 内建，用 \texttt{help cd} 查看帮助；目录以 \texttt{-} 开头时用 \texttt{cd ./-dirname}。

\subsection{\texttt{cp} —— 复制文件或目录} 语法：\texttt{cp [选项] 源 目标}

高频参数：
\begin{itemize}
\item \texttt{-r} / \texttt{-R}：递归复制目录（必需）。
\item \texttt{-p}：保留属性（时间戳、权限、所有者）。
\item \texttt{-i}：覆盖前提示。
\item \texttt{-f}：强制覆盖（危险）。
\item \texttt{-u}：增量复制（源更新或目标不存在）。
\item \texttt{-a}：归档模式（\texttt{-dpR}），适合备份。
\item \texttt{-s}：创建符号链接。
\item \texttt{-l}：创建硬链接。
\item \texttt{-v}：显示过程。
\item \texttt{-b}：覆盖前备份（加 \texttt{\string~}）。
\item \texttt{-n}：不覆盖已有文件。
\end{itemize}

进阶：
\begin{itemize}
\item \texttt{cp -au}：增量备份。
\item \texttt{cp -al}：快速创建硬链接副本（节省空间，但修改会同步）。
\item \texttt{cp --parents}：保留目录结构。
\end{itemize}

示例：
\begin{lstlisting}[language=bash]
	cp -a /data/backup /mnt/backup
	cp -i *.conf /etc/
	cp -ru ./src/* ./dst/
	cp -al /var/www /var/www_snapshot
\end{lstlisting}

注意：目标已存在目录时，\texttt{cp -r src dest} 会在 \texttt{dest} 下创建 \texttt{src} 子目录；若想合并内容，用 \texttt{cp -r src/. dest/}。

\subsection{\texttt{mv} —— 移动或重命名} 语法：\texttt{mv [选项] 源 目标}

常用参数：\texttt{-i}（提示）、\texttt{-f}（强制）、\texttt{-u}（仅更新）、\texttt{-v}、\texttt{-b}（备份）、\texttt{-n}、\texttt{-t}（指定目标目录）。

示例：
\begin{lstlisting}[language=bash]
	mv -i report.txt final_report.txt
	mv -v *.log /var/log/archive/
	mv -b existing.conf new.conf
	mv -t /backup/ *.log
\end{lstlisting}

注意：同一文件系统内移动只是改目录项，跨文件系统则复制再删除。

\subsection{\texttt{rm} —— 删除文件或目录（极度危险）} 语法：\texttt{rm [选项] 目标}

参数：\texttt{-r}（递归）、\texttt{-f}（强制）、\texttt{-i}（询问）、\texttt{-v}、\texttt{-d}（删空目录）、\texttt{--one-file-system}（防误删挂载点）。

安全实践：
\begin{itemize}
\item 设置别名 \texttt{alias rm='rm -i'}。
\item 使用 \texttt{rm -rf} 前再三确认路径。
\item 脚本中删除变量路径前务必判断非空。
\end{itemize}

示例：
\begin{lstlisting}[language=bash]
	rm -ri ./temp_folder/
	rm -f *.log
	rm -d empty_dir
\end{lstlisting}

\subsection{\texttt{mkdir} —— 创建目录} 参数：\texttt{-p}（递归创建父目录）、\texttt{-v}、\texttt{-m}（设置权限，如 \texttt{-m 755}）。

示例：
\begin{lstlisting}[language=bash]
	mkdir -p /data/project/{src,bin,log,conf}
	mkdir -p dir{1..10}
	mkdir -m 700 secret
\end{lstlisting}

\subsection{\texttt{touch} —— 创建空文件或更新时间戳} 参数：\texttt{-a}（仅访问时间）、\texttt{-m}（仅修改时间）、\texttt{-c}（不创建文件）、\texttt{-t}（指定时间，格式 \texttt{[[CC]YY]MMDDhhmm[.ss]}）、\texttt{-d}（日期字符串）、\texttt{-r}（参考文件）。

示例：
\begin{lstlisting}[language=bash]
	touch main.go util.go
	touch -m *.c
	touch -t 202308011200.00 file.txt
	touch -c existing.txt
\end{lstlisting}

\subsection{\texttt{ln} —— 创建链接} 参数：\texttt{-s}（软链接）、\texttt{-f}（覆盖）、\texttt{-i}、\texttt{-v}、\texttt{-n}（将链接名视为普通文件）、\texttt{-r}（相对路径软链接）。

软硬链接区别：软链接可跨文件系统、可指向目录，删除源失效；硬链接不可跨文件系统、不可指向目录，删除任意链接不影响数据。

示例：
\begin{lstlisting}[language=bash]
	ln -s /opt/app/bin/start.sh /usr/local/bin/start
	ln /data/file.txt /backup/file.txt
	ln -sr ./config/app.yml ./app.yml
\end{lstlisting}

注意：删除软链接用 \texttt{rm linkname}，勿加 \texttt{-r}。

\section{文本查看与内容处理}

\subsection{\texttt{cat} —— 连接并输出（小文件查看和拼接）} 参数：\texttt{-n}（编号）、\texttt{-b}（非空行编号）、\texttt{-s}（压缩空行）、\texttt{-E}（行尾显示 \texttt{\$}）、\texttt{-T}（Tab 显示为 \texttt{\^I}）、\texttt{-A}（显示所有控制字符）。

示例：
\begin{lstlisting}[language=bash]
	cat -n nginx.conf
	cat access.log.* > all_access.log
	cat -A data.txt
\end{lstlisting}

注意：大文件用 \texttt{less}。

\subsection{\texttt{less} —— 分页查看（大文件必用）} 参数：\texttt{-N}（行号）、\texttt{-S}（截断长行）、\texttt{-g}（高亮当前匹配）、\texttt{-i}（忽略大小写）、\texttt{-F}（内容少于一屏时自动退出）、\texttt{-R}（解释 ANSI 颜色）、\texttt{-X}（退出不清屏）。

快捷键：
\begin{itemize}
\item \texttt{空格/f}：下翻一页；\texttt{b}：上翻一页；\texttt{d/u}：半页。
\item \texttt{/pattern}：向下搜索；\texttt{?pattern}：向上搜索；\texttt{n/N}：下一个或上一个。
\item \texttt{g/G}：开头或末尾；\texttt{q}：退出；\texttt{F}：实时追踪（类似 \texttt{tail -f}）。
\end{itemize}

示例：
\begin{lstlisting}[language=bash]
	less -Ng application.log
	less -R docker.log
	less +F /var/log/syslog
\end{lstlisting}

\subsection{\texttt{head} —— 显示文件开头} 参数：\texttt{-n N}（前 N 行）、\texttt{-c N}（前 N 字节）、\texttt{-q}（多文件不显示文件名）、\texttt{-v}（显示文件名）。

示例：
\begin{lstlisting}[language=bash]
	head -n 20 access.log
	head -c 100 binary.bin
	head -v file1.txt file2.txt
\end{lstlisting}

\subsection{\texttt{tail} —— 显示文件末尾（看日志神器）} 参数：\texttt{-n N}（末尾 N 行）、\texttt{-f}（实时追踪）、\texttt{-F}（按文件名追踪，文件轮转也能继续）、\texttt{-c N}、\texttt{--pid=PID}（进程结束自动退出）、\texttt{--retry}。

示例：
\begin{lstlisting}[language=bash]
	tail -f /var/log/syslog
	tail -n 100 -F app.log
	tail -f /var/log/nginx/*.log
	tail -f app.log | grep --line-buffered ERROR
\end{lstlisting}

\subsection{\texttt{grep} —— 全局正则搜索（文本过滤之王）} 语法：\texttt{grep [选项] '正则模式' 文件}

关键参数：
\begin{itemize}
\item \texttt{-i}：忽略大小写；\texttt{-v}：反向匹配；\texttt{-w}：匹配整个单词。
\item \texttt{-E}：扩展正则（推荐）；\texttt{-P}：Perl 正则；\texttt{-F}：固定字符串。
\item \texttt{-r}：递归；\texttt{-l}：只列出文件名；\texttt{-o}：只输出匹配部分；\texttt{-n}：行号；\texttt{-c}：计数。
\item \texttt{-A/-B/-C}：上下文行。
\item \texttt{--include}/\texttt{--exclude}：文件过滤。
\item \texttt{--color=auto}：高亮。
\end{itemize}

示例：
\begin{lstlisting}[language=bash]
	grep -rni --include='*.conf' 'timeout' /etc/
	grep -oE '([0-9]{1,3}\.){3}[0-9]{1,3}' access.log | sort | uniq -c
	grep -C 3 'FATAL' app.log
\end{lstlisting}

性能提示：大文件用 \texttt{-F} 或 \texttt{-m} 限制输出。

\subsection{\texttt{awk} —— 文本分析与格式化（切列神器）} 语法：\texttt{awk [选项] '模式{动作}' 文件}

常用选项：\texttt{-F}（字段分隔符）、\texttt{-v}（变量赋值）、\texttt{-f}（脚本文件）。

内置变量：\texttt{\$0}（整行）、\texttt{\$1...\$NF}（列）、\texttt{NR}（行号）、\texttt{NF}（列数）、\texttt{FS/OFS}（输入与输出分隔符）、\texttt{RS/ORS}（记录分隔符）、\texttt{FILENAME}。

常见用法：
\begin{itemize}
\item 打印列：\texttt{awk '\{print \$1, \$3\}'}
\item 条件过滤：\texttt{awk '\$3 > 100'}
\item 求和：\texttt{awk '\{sum += \$1\} END \{print sum\}'}
\item 统计：\texttt{awk '\{count[\$1]++\} END \{for (i in count) print i, count[i]\}'}
\end{itemize}

示例：
\begin{lstlisting}[language=bash]
	ps aux | awk '{print $2, $11}'
	awk -F: '{print $1, $6}' /etc/passwd
	awk '{sum += $10} END {print "Total:", sum}' access.log
	awk '{ip[$1]++} END {for (i in ip) print i, ip[i]}' access.log
	awk 'length($0) > 80' file.txt
\end{lstlisting}

\subsection{\texttt{sed} —— 流编辑器（替换、插入和删除）} 语法：\texttt{sed [选项] '动作' 文件}

选项：\texttt{-i}（直接修改，建议加备份后缀）、\texttt{-n}（安静模式）、\texttt{-e}（多点编辑）、\texttt{-E}（扩展正则）、\texttt{-f}（脚本文件）。

动作：
\begin{itemize}
\item \texttt{s/旧/新/g}：替换（\texttt{g} 全局，数字指定第几个，\texttt{p} 打印）。
\item \texttt{d}：删除匹配行；\texttt{p}：打印（配合 \texttt{-n}）。
\item \texttt{a\textbackslash}：行后追加；\texttt{i\textbackslash}：行前插入；\texttt{c\textbackslash}：替换整行。
\item \texttt{=}：打印行号；\texttt{y/源/目标/}：字符转换。
\end{itemize}

地址范围：\texttt{N}（第 N 行）、\texttt{N,M}、\texttt{/pattern/}、\texttt{/start/,/end/}、\texttt{\$}（最后一行）。

示例：
\begin{lstlisting}[language=bash]
	sed 's/8080/9090/g' config.yml
	sed -i.bak 's/localhost/127.0.0.1/g' /etc/hosts
	sed -i '/^#/d; /^$/d' nginx.conf
	sed -n '/ERROR/p' app.log
	sed '3a Hello World' file.txt
	sed -E 's/[0-9]{4}-[0-9]{2}-[0-9]{2}/DATE/g' file.txt
\end{lstlisting}

\subsection{\texttt{sort} —— 排序} 参数：\texttt{-n}（数值排序）、\texttt{-r}（反向）、\texttt{-k}（指定列）、\texttt{-t}（分隔符）、\texttt{-u}（去重）、\texttt{-h}（人类可读数字排序）、\texttt{-m}（合并已排序文件）、\texttt{-o}（输出文件）、\texttt{-c}（检查是否已排序）。

示例：
\begin{lstlisting}[language=bash]
	sort -t: -k3 -n /etc/passwd
	ls -lh | sort -hrk 5
	sort -u file.txt
	sort -rnk 2 data.txt
\end{lstlisting}

\subsection{\texttt{uniq} —— 去重与统计（必须先排序）} 参数：\texttt{-c}（计数）、\texttt{-d}（仅重复行）、\texttt{-u}（仅唯一行）、\texttt{-f}（跳过字段）、\texttt{-s}（跳过字符）、\texttt{-w}（比较前 N 字符）、\texttt{-i}（忽略大小写）。

示例：
\begin{lstlisting}[language=bash]
	cat access.log | awk '{print $1}' | sort | uniq -c | sort -rn
	sort file.txt | uniq -d
	sort file.txt | uniq -u
\end{lstlisting}

\subsection{\texttt{wc} —— 统计行、词和字符} 参数：\texttt{-l}（行数）、\texttt{-w}（单词数）、\texttt{-c}（字节数）、\texttt{-m}（字符数，多字节）、\texttt{-L}（最长行长度）。

示例：
\begin{lstlisting}[language=bash]
	wc -l app.log
	wc -l *.py
	wc -w report.txt
\end{lstlisting}

\subsection{\texttt{cut} —— 按列或字符切割} 参数：\texttt{-d}（分隔符，默认 Tab）、\texttt{-f}（字段列表，如 \texttt{-f1,3}）、\texttt{-c}（字符范围，如 \texttt{-c1-5}）、\texttt{--complement}（取补集）、\texttt{-s}（不输出无分隔符的行）。

示例：
\begin{lstlisting}[language=bash]
	cut -d: -f1,6 /etc/passwd
	cut -c1-5 file.txt
	cut -d' ' -f2- file.txt
	cut -d: --complement -f3 /etc/passwd
\end{lstlisting}

\subsection{\texttt{tr} —— 字符转换、压缩、删除} 参数：\texttt{-d}（删除指定字符）、\texttt{-s}（压缩连续重复字符）、\texttt{-c}（取补集）、\texttt{-t}（截断）。

示例：
\begin{lstlisting}[language=bash]
	cat file.txt | tr 'a-z' 'A-Z'
	echo "abc123def" | tr -d '0-9'
	echo "a    b   c" | tr -s ' '
	cat file.txt | tr '\n' ' '
	tr '[:lower:]' '[:upper:]' < file.txt
\end{lstlisting}

\subsection{\texttt{xargs} —— 将标准输入转为命令行参数} 参数：\texttt{-I}（替换符，如 \texttt{-I \{\}}）、\texttt{-n}（每次传递 N 个参数）、\texttt{-p}（提示确认）、\texttt{-0}（NUL 分隔，配合 \texttt{find -print0}）、\texttt{-P}（并行进程数）、\texttt{-L}（每次 N 行）、\texttt{-t}（打印命令）。

示例：
\begin{lstlisting}[language=bash]
	find . -name "*.log" -print0 | xargs -0 rm -f
	ls *.txt | xargs -I {} mv {} {}.bak
	cat urls.txt | xargs -n 1 -P 10 wget -c
	echo "file1 file2" | xargs -n 1 touch
\end{lstlisting}

\section{查找与定位}

\subsection{\texttt{find} —— 文件系统最强搜索器} 语法：\texttt{find [路径] [表达式]}

常用测试：
\begin{itemize}
\item \texttt{-name}/\texttt{-iname}：文件名匹配。
\item \texttt{-type f/d/l}：类型。
\item \texttt{-size +100M}：大小。
\item \texttt{-mtime}/\texttt{-mmin}：修改时间。
\item \texttt{-user}/\texttt{-group}/\texttt{-perm}：用户、组和权限。
\item \texttt{-empty}：空文件或空目录。
\end{itemize}

常用操作：\texttt{-print}、\texttt{-print0}、\texttt{-exec command \{\} \textbackslash;}（或 \texttt{+}）、\texttt{-delete}、\texttt{-ls}。

逻辑组合：\texttt{-a}（与）、\texttt{-o}（或）、\texttt{!}（非）、括号。

示例：
\begin{lstlisting}[language=bash]
	find / -xdev -type f -size +1G -exec ls -lh {} \; 2>/dev/null
	find . -name '*.tmp' -delete
	find . -path ./node_modules -prune -o -name '*.js' -print
\end{lstlisting}

性能提示：使用 \texttt{-maxdepth}、\texttt{-prune} 缩小范围。

\section{进程、系统与网络}

\subsection{\texttt{ps} —— 显示进程快照} 常用组合：
\begin{itemize}
\item \texttt{ps aux}：BSD 风格，显示所有进程。
\item \texttt{ps -ef}：System V 风格。
\item \texttt{ps -e --forest}：进程树。
\item \texttt{ps -o pid,ppid,cmd}：自定义列。
\item \texttt{ps -C nginx}：按命令名查找。
\item \texttt{--sort=-\%mem}：按内存降序。
\end{itemize}

输出列（aux）：USER、PID、\%CPU、\%MEM、VSZ、RSS、TTY、STAT、START、TIME、COMMAND。

示例：
\begin{lstlisting}[language=bash]
	ps aux | grep nginx
	ps -ef --forest
	ps aux --sort=-%mem | head -6
	ps -u mysql
	ps -fp 1234
\end{lstlisting}

\subsection{\texttt{top} 与 \texttt{htop} —— 实时进程监控} 参数：\texttt{-p PID}、\texttt{-u 用户}、\texttt{-b}（批处理）、\texttt{-n 次数}、\texttt{-d 秒}、\texttt{-H}（显示线程）。

交互快捷键：\texttt{P}（CPU排序）、\texttt{M}（内存）、\texttt{N}（PID）、\texttt{T}（时间）、\texttt{k}（杀进程）、\texttt{r}（调整优先级）、\texttt{q}（退出）、\texttt{c}（完整命令）、\texttt{1}（各CPU核心）、\texttt{z}（彩色）、\texttt{f}（选择字段）。

示例：
\begin{lstlisting}[language=bash]
	top -u mysql
	top -b -n 1 > top_output.txt
\end{lstlisting}

\subsection{\texttt{kill} —— 终止进程} 常用信号：
\begin{itemize}
\item \texttt{-15} (SIGTERM)：优雅终止（默认）。
\item \texttt{-9} (SIGKILL)：强制杀死。
\item \texttt{-1} (SIGHUP)：重新加载配置。
\item \texttt{-2} (SIGINT)：类似 Ctrl+C。
\item \texttt{-19} (SIGSTOP)：暂停；\texttt{-18} (SIGCONT)：继续。
\end{itemize}

示例：
\begin{lstlisting}[language=bash]
	kill -15 12345
	kill -9 12345
	ps aux | grep java | awk '{print $2}' | xargs kill -15
	kill -1 $(cat /var/run/nginx.pid)
\end{lstlisting}

\subsection{\texttt{df} —— 查看磁盘分区空间} 参数：\texttt{-h}（人类可读）、\texttt{-H}（1000进制）、\texttt{-T}（文件系统类型）、\texttt{-i}（inode使用情况）、\texttt{--total}、\texttt{-a}（所有）、\texttt{-x}（排除类型）。

示例：
\begin{lstlisting}[language=bash]
	df -hT
	df -i
	df -h /var/log
\end{lstlisting}

\subsection{\texttt{du} —— 查看文件和目录的磁盘占用} 参数：\texttt{-s}（汇总）、\texttt{-h}、\texttt{-a}（显示所有文件）、\texttt{-c}（显示总和）、\texttt{-d}（递归深度）、\texttt{-b/-k/-m}（单位）、\texttt{--exclude}。

示例：
\begin{lstlisting}[language=bash]
	du -sh * | sort -rh
	du -h / | sort -rh | head -10
	du -sh .
	du -h -d 2 /var
\end{lstlisting}

\subsection{\texttt{free} —— 查看内存使用} 参数：\texttt{-h}、\texttt{-m/-g}、\texttt{-s}（刷新间隔）、\texttt{-t}（总计行）。

列说明：total、used、free、shared、buff/cache、available（真正可用）。

示例：
\begin{lstlisting}[language=bash]
	free -h
	free -hs 2
\end{lstlisting}

\subsection{\texttt{ss} —— 查看网络端口与连接（推荐）} 参数：\texttt{-t}（TCP）、\texttt{-u}（UDP）、\texttt{-l}（仅监听）、\texttt{-n}（数字显示）、\texttt{-p}（显示进程）、\texttt{-a}（所有）、\texttt{-r}（解析主机名）、\texttt{-s}（统计摘要）、\texttt{-4/-6}、\texttt{-o}（计时器）。

示例：
\begin{lstlisting}[language=bash]
	ss -tlnp
	ss -tan | awk '{print $1}' | sort | uniq -c
	ss -tlnp | grep :8080
\end{lstlisting}

（兼容 \texttt{netstat -tunlp}）

\subsection{\texttt{curl} —— 发送 HTTP/HTTPS 请求（接口测试）} 常用参数：
\begin{itemize}
\item \texttt{-I}：仅响应头（HEAD）。
\item \texttt{-o/-O}：保存文件。
\item \texttt{-L}：跟随重定向。
\item \texttt{-v}：详细过程。
\item \texttt{-H}：自定义请求头。
\item \texttt{-d}：POST 数据（默认 \texttt{application/x-www-form-urlencoded}）。
\item \texttt{-X}：指定方法。
\item \texttt{-F}：表单上传文件。
\item \texttt{-u}：基本认证。
\item \texttt{-x}：代理。
\item \texttt{--limit-rate}：限速。
\item \texttt{--connect-timeout}/\texttt{--max-time}：超时。
\item \texttt{-s}（静默）、\texttt{-S}（显示错误）、\texttt{-k}（忽略 SSL）、\texttt{-e}（Referer）、\texttt{-A}（User-Agent）。
\end{itemize}

示例：
\begin{lstlisting}[language=bash]
	curl -I https://example.com
	curl -X POST -H "Content-Type: application/json" -d '{"name":"test"}' http://localhost:8080/api
	curl -L -o ubuntu.iso https://releases.ubuntu.com/22.04/ubuntu-22.04-desktop-amd64.iso
	curl -F "file=@/tmp/test.txt" http://localhost/upload
	curl -u admin:secret https://api.example.com/data
	curl -v https://example.com
\end{lstlisting}

\subsection{\texttt{wget} —— 下载文件} 参数：\texttt{-c}（断点续传）、\texttt{-O}（输出文件名）、\texttt{-P}（目录）、\texttt{-q}（安静）、\texttt{-t}（重试次数）、\texttt{-i}（批量下载列表）、\texttt{--limit-rate}、\texttt{--user-agent}、\texttt{--post-data}、\texttt{-r}（递归）、\texttt{-np}（不进入上层）、\texttt{-A}（后缀过滤）。

示例：
\begin{lstlisting}[language=bash]
	wget -c https://example.com/bigfile.zip
	wget -i urls.txt
	wget --limit-rate=200k https://example.com/file.iso
\end{lstlisting}

\section{权限与打包}

\subsection{\texttt{chmod} —— 修改权限} 数字法：4=读，2=写，1=执行，组合如 755、644。 符号法：\texttt{u/g/o/a}、\texttt{+/-/=}、\texttt{r/w/x}。

参数：\texttt{-R}（递归）、\texttt{-c/-v}（显示变化）、\texttt{--reference}（参考文件）。

特殊位：SUID (4)、SGID (2)、粘滞位 (1)。

示例：
\begin{lstlisting}[language=bash]
	chmod +x deploy.sh
	chmod 644 file.txt
	find . -type d -exec chmod 755 {} \;
	find . -type f -exec chmod 644 {} \;
	chmod u+s /usr/bin/passwd
\end{lstlisting}

\subsection{\texttt{chown} —— 修改所有者和组} 参数：\texttt{-R}、\texttt{--from}（条件修改）、\texttt{-v/-c}。

示例：
\begin{lstlisting}[language=bash]
	chown www-data:www-data index.html
	chown :root file.txt
	chown -R user:admin /home/user/project
	chown --from=old new file.txt
\end{lstlisting}

\subsection{\texttt{tar} —— 打包与解包（最常用压缩工具）} 主操作：\texttt{-c}（创建）、\texttt{-x}（解包）、\texttt{-t}（列出）、\texttt{-r}（追加）、\texttt{-u}（更新）、\texttt{-d}（比较）。

辅助选项：\texttt{-v}、\texttt{-f}（归档名，必须紧跟）、\texttt{-z}（gzip）、\texttt{-j}（bzip2）、\texttt{-J}（xz）、\texttt{-C}（解压目录）、\texttt{-p}（保留权限）、\texttt{--exclude}、\texttt{--exclude-vcs}、\texttt{--strip-components}、\texttt{-P}（绝对路径，不推荐）。

示例：
\begin{lstlisting}[language=bash]
	tar -czvf project.tar.gz ./project/
	tar -xzvf project.tar.gz -C /tmp/deploy/
	tar -tzvf project.tar.gz
	tar -czvf project.tar.gz ./project --exclude='.git' --exclude='node_modules'
	tar -xzf archive.tar.gz --strip-components=1 -C /opt/
\end{lstlisting}

\section{Shell 编程核心语法（Bash）}

\subsection{变量与特殊变量}
\begin{itemize}
\item 定义：\texttt{name="Linux"}（等号两边无空格）。
\item 使用：\texttt{\$name} 或 \texttt{\$\{name\}}（推荐）。
\item 只读：\texttt{readonly name}；删除：\texttt{unset name}。
\item 特殊变量：
\end{itemize}

\begin{center}
	\begin{tabular}{ll}
		\toprule
		变量 & 含义 \\
		\midrule
		\texttt{\$0} & 脚本名称 \\
		\texttt{\$1-\$9} & 位置参数 \\
		\texttt{\$\#} & 参数个数 \\
		\texttt{\$*} & 所有参数（字符串） \\
		\texttt{\$@} & 所有参数（独立） \\
		\texttt{\$?} & 上一条命令退出状态 \\
		\texttt{\$\$} & 当前 shell PID \\
		\texttt{\$!} & 后台最后进程 PID \\
		\texttt{\$-} & 当前 shell 选项 \\
		\bottomrule
	\end{tabular}
\end{center}


示例：
\begin{lstlisting}[language=bash]
	#!/bin/bash
	echo "脚本名: $0"
	echo "第一个参数: $1"
	echo "参数个数: $#"
	shift
	echo "shift 后第一个参数: $1"
\end{lstlisting}

\subsection{条件测试} 推荐使用 \texttt{[[ 条件 ]]}。

文件测试：\texttt{-f}（普通文件）、\texttt{-d}（目录）、\texttt{-e}（存在）、\texttt{-r/-w/-x}、\texttt{-s}（非空）、\texttt{-L}（符号链接）、\texttt{-nt/-ot}（新旧比较）。

字符串测试：\texttt{-z}（空）、\texttt{-n}（非空）、\texttt{==}/\texttt{!=}、\texttt{<}/\texttt{>}（字典序）。

数值比较：\texttt{-eq}、\texttt{-ne}、\texttt{-gt}、\texttt{-lt}、\texttt{-ge}、\texttt{-le}。

逻辑：\texttt{!}、\texttt{-a/-o}（在 \texttt{[ ]} 中）、\texttt{\&\&}/\texttt{||}（在 \texttt{[[ ]] }中）。

示例：
\begin{lstlisting}[language=bash]
	if [[ -f "/etc/passwd" && -r "/etc/passwd" ]]; then
	echo "文件存在且可读"
	fi
	str="hello"
	if [[ -n $str && $str == "hello" ]]; then
	echo "匹配"
	fi
	num=10
	if (( num > 5 )); then
	echo "大于5"
	fi
\end{lstlisting}

\subsection{流程控制} if 语句：
\begin{lstlisting}[language=bash]
	if 条件; then
	命令
	elif 条件; then
	命令
	else
	命令
	fi
\end{lstlisting}

for 循环：
\begin{lstlisting}[language=bash]
	for i in 1 2 3 4 5; do echo $i; done
	for i in {1..10}; do echo $i; done
	for ((i=1; i<=10; i++)); do echo $i; done
	for file in *.txt; do echo "Processing $file"; done
\end{lstlisting}

while 与 until：
\begin{lstlisting}[language=bash]
	while 条件; do 命令; done
	until 条件; do 命令; done
\end{lstlisting}

case 分支：
\begin{lstlisting}[language=bash]
	case $1 in
	start) echo "启动..." ;;
	stop)  echo "停止..." ;;
	*)     echo "Usage: $0 {start|stop}"; exit 1 ;;
	esac
\end{lstlisting}

循环控制：\texttt{break}、\texttt{continue}。

\subsection{函数与算术运算} 函数：
\begin{lstlisting}[language=bash]
	function add() {
		local a=$1
		local b=$2
		echo $((a + b))
		return 0
	}
	result=$(add 3 5)
	echo "结果: $result"
\end{lstlisting}

算术运算：\texttt{\$((表达式))}、\texttt{let}、\texttt{expr}。
\begin{lstlisting}[language=bash]
	a=10; b=20
	sum=$((a + b))
	((a++))
	let "b+=5"
	c=$(expr $a + $b)
\end{lstlisting}

\subsection{数组}
\begin{lstlisting}[language=bash]
	arr=(apple banana cherry)
	echo ${arr[0]}
	echo ${arr[@]}
	echo ${#arr[@]}
	arr[3]="date"
	unset arr[1]
\end{lstlisting}

\subsection{调试技巧} 在脚本开头添加：
\begin{lstlisting}[language=bash]
	set -euxo pipefail
\end{lstlisting}
含义：\texttt{-e}（出错退出）、\texttt{-u}（未定义变量报错）、\texttt{-x}（打印命令）、\texttt{-o pipefail}（管道失败即失败）。

也可使用 \texttt{trap} 捕获信号。

\section{高级黄金组合拳（拿来即用）}
\begin{itemize}
\item 批量杀死 Java 进程：\\ \texttt{ps aux | grep 'java' | awk '\{print \$2\}' | xargs kill -9}
\item 查看某端口占用并杀进程：\\ \texttt{kill -9 \$(lsof -t -i:8080)} \\ \texttt{ss -tlnp | grep :8080 | awk '\{print \$7\}' | cut -d/ -f1 | xargs kill -9}
\item 统计日志中状态码分布：\\ \texttt{cat access.log | awk '\{print \$9\}' | sort | uniq -c | sort -rn}
\item 清理 7 天前大于 100M 的日志：\\ \texttt{find /var/log -name "*.log" -mtime +7 -size +100M -exec rm -f \{\} \textbackslash;}
\item 实时监控日志错误数：\\ \texttt{tail -f app.log | grep --line-buffered ERROR | wc -l}
\item 不落地传输整个文件夹到远程：\\ \texttt{tar -czvf - ./folder | ssh user@ip "tar -xzvf - -C /target"}
\item 提取所有 IP 并去重：\\ \texttt{grep -oE '[0-9]\{1,3\}\textbackslash.[0-9]\{1,3\}\textbackslash.[0-9]\{1,3\}\textbackslash.[0-9]\{1,3\}' access.log | sort | uniq}
\item 查找当前目录下最大的 5 个文件：\\ \texttt{find . -type f -exec ls -s \{\} \textbackslash; | sort -rn | head -5 | awk '\{print \$2\}'} \\ \texttt{du -sh * | sort -rh | head -5}
\end{itemize}

\section{GCC 使用简介}

GCC 是 GNU C 编译器。以 \texttt{main.c} 为例，编译过程可以分为四步：

\begin{itemize}
\item \texttt{gcc -E main.c}：预处理，展开头文件与宏；
\item \texttt{gcc -S main.c}：生成汇编 \texttt{main.s}；
\item \texttt{gcc -c main.c}：汇编并生成目标文件 \texttt{main.o}；
\item \texttt{gcc main.o -o app}：链接生成可执行文件 \texttt{app}。
\end{itemize}

常用选项：

\begin{table}[htbp]
\centering \small \caption{GCC 常用选项} \label{tab:gcc-option}
    \begin{tabular}{lll}
        \toprule
        选项 & 示例 & 作用 \\
        \midrule
        \texttt{-o} & \texttt{gcc main.c -o app} & 指定输出文件名 \\
        \texttt{-Wall -Wextra} & 常与 \texttt{-g} 同用 & 打开警告 \\
        \texttt{-g} & \texttt{-g} & 生成调试信息 \\
        \texttt{-O0/-O2/-Os} & \texttt{-O2} & 优化等级 \\
        \texttt{-I} & \texttt{-I./inc} & 添加头文件搜索路径 \\
        \texttt{-L/-l} & \texttt{-L./lib -lm} & 库路径与链接库 \\
        \texttt{-D} & \texttt{-DDEBUG=1} & 定义宏 \\
        \texttt{-std} & \texttt{-std=c11} & 指定 C 标准 \\
        \texttt{-c} & \texttt{gcc -c a.c} & 只生成目标文件 \\
        \bottomrule
    \end{tabular}
\end{table}

多文件编译：

\begin{lstlisting}
gcc -Wall -Wextra -g main.c uart.c pid.c -o app
\end{lstlisting}

交叉编译时把 \texttt{gcc} 换成目标工具链，例如：

\begin{lstlisting}
arm-none-eabi-gcc -mcpu=cortex-m4 -mthumb main.c -o app.elf
aarch64-linux-gnu-gcc main.c -o app_arm64
\end{lstlisting}

% ---------------------------------------------------------------------
\section{Makefile 使用简介}

手动输入 GCC 命令容易漏文件，Makefile 用“目标、依赖和命令”描述构建规则：

\begin{lstlisting}
CC      = gcc
CFLAGS  = -Wall -Wextra -g
TARGET  = app
SRCS    = main.c uart.c pid.c
OBJS    = $(SRCS:.c=.o)

$(TARGET): $(OBJS)
	$(CC) $(CFLAGS) $^ -o $@

%.o: %.c
	$(CC) $(CFLAGS) -c $< -o $@

clean:
	rm -f $(OBJS) $(TARGET)

.PHONY: clean
\end{lstlisting}

\begin{itemize}
\item \texttt{\$@}：目标名；
\item \texttt{\$<}：第一个依赖；
\item \texttt{\${\string^}}：全部依赖；
\item \texttt{.PHONY}：告诉 make \texttt{clean} 不是真实文件。
\end{itemize}

使用时只需 \texttt{make} 构建、\texttt{make clean} 清理。Makefile 适合中小工程；工程文件变多、需要跨平台时更推荐 CMake。

% ---------------------------------------------------------------------
\section{CMake 使用简介}

CMake 先生成构建系统，再调用 GCC、Make 或 Ninja。最小工程：

\begin{lstlisting}
cmake_minimum_required(VERSION 3.16)
project(demo C)

set(CMAKE_C_STANDARD 11)
set(CMAKE_C_STANDARD_REQUIRED ON)

add_executable(app
    main.c
    uart.c
    pid.c
)

target_include_directories(app PRIVATE inc)
target_compile_options(app PRIVATE -Wall -Wextra)
\end{lstlisting}

使用方式：

\begin{lstlisting}
cmake -S . -B build
cmake --build build
./build/app
\end{lstlisting}

\begin{itemize}
\item \texttt{cmake -S . -B build}：指定源码目录与构建目录；
\item \texttt{add\_executable}：生成可执行文件；
\item \texttt{add\_library}：生成静态库和动态库；
\item \texttt{target\_include\_directories}：添加头文件路径；
\item 交叉编译时通过工具链文件设置编译器前缀。
\end{itemize}

CMake 的价值是“一处写 CMakeLists.txt，Windows、Linux 与单片机交叉编译都能复用”。

% ---------------------------------------------------------------------
\section{后续转行——Linux 驱动与应用开发}

对从单片机转过来的开发者，最自然的第一步不是看内核源码，而是先把 RK3588 与 K230 开发板跑起来，写一个“用户空间点亮 LED”的程序，再逐步下沉到驱动层。本章的目标是“快速上手”，但 Linux 体系庞大，单靠一章命令速查远远不够。以下书目按学习路径分层整理，供读者根据自身阶段选读。

\subsection*{第一阶段：命令行与系统基础}
\begin{itemize}
	\item 《鸟哥的 Linux 私房菜》（基础学习篇） —— 鸟哥著，人民邮电出版社。Linux 入门经典，从系统安装、文件系统到 Shell 脚本，讲解细致，适合零基础读者建立系统管理直觉。
	\item 《Linux 命令行与 Shell 脚本编程大全》（第 3 版） —— Richard Blum 等著，人民邮电出版社。命令参考与脚本编程并重，适合作为案头工具书随时查阅。
	\item 《Linux 就该这么学》 —— 刘遄著，人民邮电出版社。以 RHEL/CentOS 为蓝本，结合大量实例讲解常见服务配置，适合动手型读者。
\end{itemize}

\subsection*{第二阶段：Linux 应用层与网络编程}
\begin{itemize}
	\item 《UNIX 环境高级编程》（第 3 版，APUE） —— W. Richard Stevens 著，人民邮电出版社。系统编程的“圣经”，涵盖文件 I/O、进程控制、信号、线程、进程间通信等，C 语言开发者必读。
	\item 《Linux C 与 C++ 一线开发实践》 —— 朱文伟等著，清华大学出版社。贴近企业一线开发实践，涵盖 Linux 文件编程、多进程、进程间通信及 Web 编程，适合从单片机转向 Linux 应用开发的读者。
	\item 《Linux 网络编程》 —— 宋敬彬编著，清华大学出版社（“Linux 典藏大系”）。全面覆盖套接字编程、高级 I/O 函数、服务器模型等，配有教学视频与习题。
\end{itemize}

\subsection*{第三阶段：嵌入式 Linux 系统与驱动开发}
\begin{itemize}
	\item 《嵌入式 Linux 应用开发完全手册》 —— 韦东山著，人民邮电出版社。以项目驱动，从 U-Boot 移植、内核裁剪到驱动编写，偏重工程实操，是嵌入式 Linux 入门的首选。
	\item 《Linux 设备驱动程序》（第 3 版，LDD3） —— Jonathan Corbet 等著，O'Reilly。驱动开发的经典教材，以实践为导向讲解字符设备、块设备、网络接口驱动，覆盖中断、DMA、内存分配等核心主题[reference:5]。虽然内核版本较老，但设计思想至今仍有效。
	\item 《Linux 设备驱动开发详解》（第 2 版） —— 宋宝华著，人民邮电出版社。结合 Linux 内核源码，系统讲解驱动框架与实现原理，豆瓣评分 8.2，适合深入阅读。
	\item 《Linux 设备驱动开发》（第 2 版） —— John Madieu 著，陈莉君译，人民邮电出版社。基于最新内核版本，涵盖 IIO、RTC、DMA 等子系统，理论结合大量源码示例。
\end{itemize}

\subsection*{第四阶段：Linux 内核原理（进阶）}
\begin{itemize}
	\item 《深入理解 Linux 内核》（第 3 版） —— Daniel P. Bovet 等著，中国电力出版社。全面剖析内核核心子系统——进程调度、内存管理、文件系统、网络栈等，适合有一定基础后深入研读。
	\item 《Linux 内核设计与实现》（原书第 3 版，典藏版） —— Robert Love 著，机械工业出版社。篇幅精炼，行文清晰，适合快速建立内核整体框架认知。
	\item 《图解 Linux 内核：基于 6.x》 —— 姜亚华著，机械工业出版社。2024 年出版，基于 Linux 6.x 内核，图文并茂，适合希望结合最新内核代码学习的读者。
\end{itemize}

\subsection*{选读建议}
从单片机转向嵌入式 Linux，不必一上来就啃内核。务实的路径是：先跑通开发板、写用户空间程序（第二阶段）\(\rightarrow\) 再学 U-Boot 与内核移植（第三阶段）\(\rightarrow\) 最后按需深入驱动或内核（第四阶段）。第一阶段的命令和 Shell 脚本贯穿始终，建议边用边查，不必一次性背完。
% ---------------------------------------------------------------------
\section{实验报告：Linux 开发环境搭建与基础命令验证}

\subsection{实验基本信息}

Linux 发行版：\LabFillLine[3.5cm] \quad 内核版本：\LabFillLine[3.5cm]

虚拟机软件：\LabFillLine[3.5cm] \quad 分配内存/CPU：\LabFillLine[3.5cm]

GCC 版本：\LabFillLine[3.5cm] \quad CMake 版本：\LabFillLine[3.5cm]

\subsection{实验目的}
\begin{enumerate}
\item 掌握 Linux 虚拟机的安装与基本配置方法；
\item 熟悉 Linux 常用命令（文件操作、文本处理、进程管理、网络工具）；
\item 掌握 GCC 单文件与多文件编译流程；
\item 理解 Makefile 自动构建规则；
\item 掌握 CMake 工程构建方法。
\end{enumerate}

\subsection{实验原理}
Linux 是嵌入式开发的重要平台，现代工具链（GCC、CMake、Git 等）均原生支持 Linux 环境。虚拟机技术允许在 Windows 主机上运行完整的 Linux 系统，无需双系统或额外硬件。GCC 将 C 源码依次经过预处理、编译、汇编、链接四个阶段生成可执行文件。Makefile 通过“目标-依赖-命令”规则实现增量编译，CMake 则进一步抽象构建过程，生成平台无关的构建配置。

\subsection{硬件连接}

\LabPasteBox{Linux 开发环境硬件连接}{粘贴虚拟机网络配置截图（NAT/桥接模式）和主机与虚拟机通信测试（ping）截图。}{fig:linux-hardware}

\subsection{关键代码}

\LabPasteBox{Linux 基础命令与构建脚本}{粘贴 Shell 测试命令、GCC 编译命令、Makefile 和 CMakeLists.txt 核心内容。}{fig:linux-code}

\subsection{实验步骤}
\begin{enumerate}
\item 下载 Ubuntu 22.04/24.04 ISO 镜像，在 VMware 或 VirtualBox 中创建虚拟机，分配 4GB 以上内存和 60GB 以上磁盘。
\item 安装系统，配置用户和网络（NAT 模式确保能联网），执行 \texttt{sudo apt update \&\& sudo apt upgrade -y} 更新系统。
\item 安装开发工具：\texttt{sudo apt install -y build-essential cmake git openssh-server}。
\item 练习文件操作命令：\texttt{ls}、\texttt{cd}、\texttt{cp}、\texttt{mv}、\texttt{rm}、\texttt{mkdir}、\texttt{touch}、\texttt{ln}。
\item 练习文本查看与处理命令：\texttt{cat}、\texttt{less}、\texttt{head}、\texttt{tail}、\texttt{grep}、\texttt{awk}、\texttt{sed}、\texttt{sort}、\texttt{uniq}、\texttt{wc}、\texttt{cut}、\texttt{tr}、\texttt{xargs}。
\item 练习查找命令：\texttt{find}，以及进程与网络命令：\texttt{ps}、\texttt{top}、\texttt{kill}、\texttt{df}、\texttt{du}、\texttt{free}、\texttt{ss}、\texttt{curl}、\texttt{wget}。
\item 编写一个简单的 C 程序（\texttt{hello.c}），分别使用 GCC 单步（\texttt{-E}、\texttt{-S}、\texttt{-c}、\texttt{-o}）和一步编译，生成可执行文件并运行。
\item 编写 Makefile，实现 \texttt{hello.c} 的自动构建和清理，执行 \texttt{make} 和 \texttt{make clean} 验证。
\item 编写 \texttt{CMakeLists.txt}，使用 \texttt{cmake -S . -B build} 和 \texttt{cmake --build build} 构建工程。
\item 启用 SSH 服务，在主机上用 VS Code Remote-SSH 连接虚拟机，验证远程开发环境。
\end{enumerate}

\subsection{实验现象与数据记录}

\begin{table}[htbp]
\centering
\caption{Linux 命令与构建工具验证记录}
\label{tab:linux-record}
\begin{tabularx}{\textwidth}{p{2.8cm}p{3.5cm}p{3.5cm}X}
\toprule
\textbf{测试项目} & \textbf{输入命令/操作} & \textbf{预期结果} & \textbf{实际结果} \\
\midrule
文件操作 & \texttt{ls -la \string~} & 显示家目录所有文件 & \LabFillLine[2.8cm] \\
文本搜索 & \texttt{grep -rn "main" .} & 显示含 main 的行 & \LabFillLine[2.8cm] \\
GCC 编译 & \texttt{gcc hello.c -o hello} & 生成可执行文件 & \LabFillLine[2.8cm] \\
Makefile & \texttt{make} & 自动编译生成目标 & \LabFillLine[2.8cm] \\
CMake & \texttt{cmake --build build} & 构建成功 & \LabFillLine[2.8cm] \\
\bottomrule
\end{tabularx}
\end{table}

\LabPasteBox{Linux 实验运行结果}{粘贴终端执行 \texttt{uname -a}、GCC 编译、Make 构建和 CMake 构建的完整输出截图。}{fig:linux-result}

\subsection{结果分析}

虚拟机网络配置（NAT/桥接模式）的选择依据：\LabFillLine[6cm]

GCC 预处理、编译、汇编、链接四个阶段分别生成什么文件：\LabFillLine[6cm]

Makefile 中 \texttt{\$@}、\texttt{\$<}、\texttt{\${\string^}} 分别代表什么：\LabFillLine[6cm]

CMake 相比 Makefile 在跨平台构建中的优势：\LabFillLine[6cm]

\subsection{实验总结}

\noindent\LabFillLine[\linewidth]\par
\noindent\LabFillLine[\linewidth]\par
\noindent\LabFillLine[\linewidth]

% ---------------------------------------------------------------------
\section{小结}

Linux 对嵌入式工程师不是“另一个操作系统”，而是“新的工作台”。本章只覆盖了最高频的命令、GCC、Makefile 与 CMake；把它们练熟后，再学交叉编译、设备树和内核驱动会顺利得多。

% ---------------------------------------------------------------------
\begin{problemset}
\item 用一条命令组合完成：“递归查找 \texttt{src/} 下所有包含 \texttt{error} 的 .c 文件并显示行号”。
\item 说明 \texttt{ps aux | grep python}、\texttt{awk '\{print \$2\}'} 与 \texttt{xargs kill} 组合的含义，并说明为什么需要 \texttt{xargs}。
\item 写一个 Bash 脚本：遍历 \texttt{/tmp/logs} 下所有 \texttt{.log} 文件，删除其中超过 7 天且大于 1M 的文件，并打印处理结果。
\item Makefile 中的“目标、依赖和命令”三者关系是什么？为什么修改头文件后，依赖写全的工程能自动重新编译？
\item 用 CMake 构建一个包含 \texttt{src} 与 \texttt{include} 两个目录的最小 C 工程，并解释 \texttt{cmake -S . -B build} 的作用。
\item GCC 的 \texttt{-Wall -Wextra -g} 分别解决了什么问题？调试器需要的调试信息由哪个选项生成？
\end{problemset}
